\documentclass[11pt]{letter}

\usepackage[margin=0.75in]{geometry}
\usepackage{setspace}
\usepackage{hyperref}
\usepackage{xcolor}
\usepackage{enumitem}
\usepackage{parskip}

\hypersetup{
    colorlinks=true,
    urlcolor=black,
    linkcolor=black
}

\setlength{\parindent}{0pt}
\setstretch{1.05}
\date{}
\begin{document}

\begin{letter}{}

\opening{Dear Hiring Committee,}

I saw the job post on \textbf{LinkedIn} and thought I'd be a good fit for it.

I am a Master's student in Computer Science at the University of Cincinnati, graduating December 2026, looking for a role as a \textbf{Software Engineer, Data and Automation} at \textbf{PNC}. With hands on experience building Python automation, cloud infrastructure, and CI/CD release pipelines on AWS, I am excited about the opportunity to contribute to \textbf{PNC} and support the Data and Automation organization's mission of scalable, risk-aware software solutions.

\textbf{Here are reasons that make me a great fit for the role:}

\begin{itemize}[leftmargin=*, itemsep=0.45em]
    \item \textbf{Technical design and development of complex solutions:} At Siemens Industry Software, I designed and developed Teamcenter X infrastructure automation with Python and Boto3 on AWS, replacing legacy Ruby scripts and cutting deployment time by 40\%, while executing 44 production environment upgrades in 4 months.

    \item \textbf{Collaboration with product owners, architects, and engineers:} I worked with architects and engineers to design an Ansible based Blue-Green deployment solution that reduced production upgrade downtime by 50\%, and migrated 40+ production environments to Kubernetes using ArgoCD, Helm, and GitLab CI/CD in alignment with enterprise standards.

    \item \textbf{Agile delivery, CI/CD, and automated testing:} I built CI/CD release pipelines with GitLab and GitHub integrated workflows, and as Quality Coach and SonarQube Specialist I drove automated testing adoption that raised code coverage from 35\% to 78\% across an Agile Release Train.

    \item \textbf{Python and AWS application development:} My graduate research applies Python to a memory-augmented LLM pipeline for Android bug reproduction, and my NYC Mobility Forecasting project used Python and PySpark to build a demand forecasting system across 263 zones on a 2.3M-row dataset, cutting RMSE by 43\% over baseline.

    \item \textbf{Technical documentation and process improvement:} I authored design artifacts and onboarding documentation while mentoring 5 engineers, and consistently focused on process improvements that reduced deployment and downtime metrics across production systems.
\end{itemize}

I am available to work with you after my graduation in December 2026. As a Master's student at the University of Cincinnati and someone who values professionalism and growth, I am enthusiastic about contributing to PNC's Data and Automation team.

Please find my CV attached below. I look forward to hearing from you. I would be excited to bring my Python, AWS, and automation experience to a team building scalable solutions for PNC's customers.

{\setlength{\parskip}{0pt}
Regards,\\
Tanmay Bhuskute,\\
tanmay.v.bhuskute@gmail.com $\mid$ +1 (732) 858-3027,\\
Cincinnati, OH.
}

\end{letter}

\end{document}
